% =============================================================================
\bloco{3}{Created: the Prova~01 economy, extended}{Kurlat \S1.1 \,\textbullet\, 35 points}

\questao{(12 points) \emph{Inventories and data sales.} The economy of Prova~01 in a new year. All values in reais;
nothing else happens.
\begin{itemize}[leftmargin=1.4em]
\item Firm A uses no intermediate inputs. It produces inputs worth 800 and pays wages of 500. It sells inputs worth 600 to B and
keeps the remaining 200 in its warehouse.
\item Firm B uses all 600 of the inputs it bought, pays wages of 200, buys advertising from N for 100 and a market-research
dataset from N for 60. It sells final goods to households for 900; final goods worth 150 at market prices are still unsold in its
warehouse at the end of the year.
\item Social network N is free for its users. It sells the advertising (100) and the dataset (60) to B and another dataset to a
firm abroad for 40. It pays wages of 120 and has no other costs.
\end{itemize}}
\noindent\textbf{a)} (7 points) Compute GDP by the product, income and expenditure approaches in one table (by firm for product and
income, by component for expenditure) and check that they agree.\\
\textbf{b)} (3 points) In words: why do the 200 of stored inputs and the 150 of unsold goods count in this year's GDP, and
why does N's data enter the expenditure column once and not the other time?\\
\textbf{c)} (2 points) Next year B sells its 150 of stored goods to households for 150, and nothing else happens. Give the
three approaches for that year and explain the result in one sentence.
\resolespaco
\begin{sol}
\begin{notacao}
\item A, B, N: the input firm, the final-goods firm and the social network. $Q_j$, $Z_j$, $VA_j$, $W_j$, $\Pi_j$ for
$j\in\{A,B,N\}$, and the identities \eqref{q1:G1}--\eqref{q1:G4}: as in Question~1.
\item $C$, $X\ge0$: consumption and exports; $G=M=0$. $I$: investment, here the change in inventories (negative when they are
drawn down).
\end{notacao}

\passo{a) value added}{Gross output is sales plus additions to inventory. A sells 600 and stores 200:}
\begin{align}
Q_A &= (600)+(200) \nonumber\\
&= 800. \label{q7:QA}
\end{align}
A uses no inputs, so by \eqref{q1:G2} and \eqref{q7:QA}:
\begin{align*}
VA_A &= (800)-(0)\\
&= 800 .
\end{align*}
B sells 900 and stores 150:
\begin{align}
Q_B &= (900)+(150) \nonumber\\
&= 1{,}050. \label{q7:QB}
\end{align}
B's intermediate inputs are A's inputs, the advertising and the dataset:
\begin{align}
Z_B &= (600)+(100)+(60) \nonumber\\
&= 760. \label{q7:ZB}
\end{align}
Using \eqref{q1:G2} with \eqref{q7:QB} and \eqref{q7:ZB}:
\begin{align}
VA_B &= (1{,}050)-(760) \nonumber\\
&= 290. \label{q7:vaB}
\end{align}
N sells the advertising and the two datasets, with no inputs:
\begin{align}
Q_N=VA_N &= (100)+(60)+(40) \nonumber\\
&= 200. \label{q7:vaN}
\end{align}
Summing $VA_A=800$, \eqref{q7:vaB} and \eqref{q7:vaN} with \eqref{q1:G1}:
\begin{align*}
\text{GDP} &= (800)+(290)+(200)\\
&= 1{,}290 .
\end{align*}
\passo{a) income}{By \eqref{q1:G3}, each firm's profit is its value added minus its wages:}
\begin{align}
\Pi_A &= (800)-(500) \nonumber\\
&= 300, \label{q7:PiA}\\
\Pi_B &= (1{,}050)-(760)-(200) \nonumber\\
&= 90, \label{q7:PiB}\\
\Pi_N &= (200)-(120) \nonumber\\
&= 80. \label{q7:PiN}
\end{align}
Adding up the wages and the profits \eqref{q7:PiA}--\eqref{q7:PiN}:
\begin{align}
\textstyle\sum_j W_j &= (500)+(200)+(120) \nonumber\\
&= 820, \label{q7:W}\\
\textstyle\sum_j \Pi_j &= (300)+(90)+(80) \nonumber\\
&= 470. \label{q7:Pi}
\end{align}
By \eqref{q1:G3}, the sum of \eqref{q7:W} and \eqref{q7:Pi}:
\begin{align*}
\text{GDP} &= (820)+(470)\\
&= 1{,}290 .
\end{align*}
\passo{a) expenditure}{Investment is the change in inventories, A's 200 plus B's 150:}
\begin{align}
I &= (200)+(150) \nonumber\\
&= 350. \label{q7:I}
\end{align}
With $C=900$, $X=40$ and $G=M=0$, \eqref{q1:G4} and \eqref{q7:I} give
\begin{align*}
\text{GDP} &= (900)+(350)+(40)\\
&= 1{,}290 .
\end{align*}
\begin{center}\footnotesize\begin{tabular}{@{}lr@{\qquad}lr@{\qquad}lr@{}}\toprule
\multicolumn{2}{@{}c}{Production (value added)} & \multicolumn{2}{c}{Income} & \multicolumn{2}{c@{}}{Expenditure}\\\midrule
A: $800-0$ & 800 & A: wages 500 $+$ profit 300 & 800 & Consumption & 900\\
B: $1{,}050-760$ & 290 & B: wages 200 $+$ profit 90 & 290 & Investment (inventories) & 350\\
N: $200-0$ & 200 & N: wages 120 $+$ profit 80 & 200 & Exports & 40\\\midrule
GDP & 1{,}290 & GDP & 1{,}290 & GDP & 1{,}290\\\bottomrule\end{tabular}\end{center}
\passo{b) interpretation}{GDP counts what is \emph{produced} this year, sold or not. Goods produced and kept are an investment in
inventories, valued in the producer's output and in its profit; that is how the product and income columns pick them up. The
dataset sold to B is an input B uses up in production --- intermediate, already inside the price of B's goods --- so it has no
final-use entry. The dataset sold abroad leaves the country: it is a final use by a non-resident, an export. The service users get
for free has no price and stays out, as in Prova~01.}
\passo{c) next year}{Production: B's output is its sales minus the goods drawn from inventory, so by \eqref{q1:G2}:}
\begin{align*}
VA_B=Q_B &= (150)-(150)\\
&= 0 .
\end{align*}
Income: B's revenue 150 against the 150 value of goods leaving the warehouse; by \eqref{q1:G3}:
\begin{align*}
\Pi_B &= (150)-(150)\\
&= 0 .
\end{align*}
Expenditure: $C=+150$ and $I=-150$ (inventories fall); by \eqref{q1:G4}:
\begin{align*}
\text{GDP} &= (150)+(-150)\\
&= 0 .
\end{align*}
GDP $=0$: the goods were counted the year they were made; selling them later moves them from $I$ to $C$ and creates no new
production.
\intuicao{The 160 of N's sales to B and the 600 of A's sales to B appear in the product column only as deductions from B's output;
the final-use column sees only what is left at the end of the chain: consumption, inventories and exports.}
\refbib{Kurlat (2020), \S1.1, Ex.~1.2 (value added), p.~17; Ex.~1.3 (inventories as investment, negative when drawn down), p.~17;
Ex.~1.6 (exports), p.~19. Same treatment of inventories in Williamson (2014), solutions manual, ch.~2, problems 1, 3 and 6
(``additions to inventory are treated as purchasing one's own output'').}
\rubrica{
\item a) 7 points: each approach 2, agreement checked 1. Only one approach: max.\ 2.5 (the Prova~01 Q1(b) error).
\item Inventories left out (GDP $=940$), or the data sold to B added to $C$: max.\ 3 for a).
\item b) 3 points: inventories 1.5, the two data sales 1.5. c) 2 points; GDP of 150 next year: 0 for c).
\item Arithmetic slip with the right layout: $-0.5$, once.
}
\end{sol}

% -----------------------------------------------------------------------------
\questao{(12 points) \emph{Imports, a software sale, dividends and the government.} Another year of the Prova~01 economy.
All values in reais; nothing else happens.
\begin{itemize}[leftmargin=1.4em]
\item Firm A imports raw materials for 100, turns them into inputs and sells the inputs to B for 600. It pays wages of 350.
\item Firm B uses the inputs, buys advertising from N for 100, pays wages of 250 and sells final goods for 1{,}000: 850 to
domestic households and 150 to buyers abroad.
\item Network N is free for its users. It sells the advertising (100) and sells licences of a photo-editing app to domestic
households for 80. It pays wages of 150, of which 100 go to a team that builds a recommendation engine N will use itself for
several years. It has no other costs. Out of its profit, N pays dividends of 40; half of its shareholders live abroad.
\item The government hires content moderators for wages of 30 (their service is free to the public), pays unemployment
benefits of 25 and collects taxes of 55 from households.
\end{itemize}}
\noindent\textbf{a)} (6 points) Compute GDP by the three approaches in one table and check that they agree.\\
\textbf{b)} (2 points) Compute gross national income (GNI) and explain why it differs from GDP.\\
\textbf{c)} (2 points) In words: classify the app licences, the recommendation engine, the dividends and the unemployment
benefits, each in one sentence.\\
\textbf{d)} (2 points) Prova~01 Q1(c) had GDP $=1{,}100$. Account for the difference between that number and your answer
to (a), item by item.
\resolespaco
\begin{sol}
\begin{notacao}
\item A, B, N, and Gov (the government as a producer). $Q_j$, $Z_j$, $VA_j$, $W_j$, $\Pi_j$ (operating surplus for N) and
the identities \eqref{q1:G1}--\eqref{q1:G4}: as in Question~1.
\item $C$, $I$, $G$, $X$, $M\ge0$: consumption, investment, government purchases, exports, imports.
\item GNI: gross national income, the income of residents. $NFI$: net primary income from abroad (negative when more is paid
to non-residents than received from them).
\end{notacao}

\passo{a) value added}{By \eqref{q1:G2}, output minus intermediate inputs. A's input is the imported raw material; B's are A's
inputs and the advertising:}
\begin{align}
VA_A &= (600)-(100) \nonumber\\
&= 500, \label{q8:vaA}\\
VA_B &= (1{,}000)-(600)-(100) \nonumber\\
&= 300. \label{q8:vaB}
\end{align}
N has no inputs; its output is the ads, the app and the own-account software, valued at its cost (the team's wages):
\begin{align}
VA_N &= (100)+(80)+(100) \nonumber\\
&= 280. \label{q8:vaN}
\end{align}
The government's output is valued at cost, $VA_{Gov}=30$. Summing with \eqref{q1:G1}:
\begin{align*}
\text{GDP} &= (500)+(300)+(280)+(30)\\
&= 1{,}110 .
\end{align*}
\passo{a) income}{By \eqref{q1:G3}, profit is value added minus wages, using \eqref{q8:vaA}--\eqref{q8:vaN}:}
\begin{align*}
\Pi_A &= (600)-(100)-(350)\\
&= 150,\\
\Pi_B &= (1{,}000)-(600)-(100)-(250)\\
&= 50,\\
\Pi_N &= (280)-(150)\\
&= 130.
\end{align*}
The government pays wages of 30 and earns no surplus. Adding up wages and profits:
\begin{align}
\textstyle\sum_j W_j &= (350)+(250)+(150)+(30) \nonumber\\
&= 780, \label{q8:W}\\
\textstyle\sum_j \Pi_j &= (150)+(50)+(130) \nonumber\\
&= 330. \label{q8:Pi}
\end{align}
By \eqref{q1:G3}, the sum of \eqref{q8:W} and \eqref{q8:Pi}:
\begin{align}
\text{GDP} &= (780)+(330) \nonumber\\
&= 1{,}110. \label{q8:Y}
\end{align}
\passo{a) expenditure}{Consumption is B's domestic sales plus the app licences:}
\begin{align}
C &= (850)+(80) \nonumber\\
&= 930. \label{q8:C}
\end{align}
With $I=100$ (the recommendation engine), $G=30$, $X=150$ and $M=100$, \eqref{q1:G4} and \eqref{q8:C} give
\begin{align*}
\text{GDP} &= (930)+(100)+(30)+(150)-(100)\\
&= 1{,}110 .
\end{align*}
\begin{center}\footnotesize\begin{tabular}{@{}lr@{\qquad}lr@{\qquad}lr@{}}\toprule
\multicolumn{2}{@{}c}{Production (value added)} & \multicolumn{2}{c}{Income} & \multicolumn{2}{c@{}}{Expenditure}\\\midrule
A: $600-100$ & 500 & A: wages 350 $+$ profit 150 & 500 & Consumption $850+80$ & 930\\
B: $1{,}000-600-100$ & 300 & B: wages 250 $+$ profit 50 & 300 & Investment (software) & 100\\
N: $100+80+100$ & 280 & N: wages 150 $+$ surplus 130 & 280 & Government & 30\\
Government (at cost) & 30 & Government: wages 30 & 30 & Exports & 150\\
 & & & & Imports & $-100$\\\midrule
GDP & 1{,}110 & GDP & 1{,}110 & GDP & 1{,}110\\\bottomrule\end{tabular}\end{center}
\passo{b) GNI}{Half of the dividends is income paid to non-residents out of domestic production, and nothing is received from
abroad:}
\begin{align*}
NFI &= -\frac{(40)}{2}\\
&= -20 .
\end{align*}
GNI adds net primary income from abroad to GDP:
\begin{equation}
\text{GNI}=\text{GDP}+NFI. \label{q8:gni}
\end{equation}
Substitute \eqref{q8:Y} and $NFI=-20$ into \eqref{q8:gni}:
\begin{align*}
\text{GNI} &= (1{,}110)+(-20)\\
&= 1{,}090 .
\end{align*}
GDP measures production inside the country; GNI measures the income of its residents.
\passo{c) classification}{App licences: a service sold to households, final consumption ($C$). Recommendation engine: own-account
software, gross fixed capital formation ($I$), valued at cost. Dividends: a distribution of N's profit already counted in the income
column; they change who receives the income, not GDP (only GNI, for the part paid abroad). Unemployment benefits: a transfer
financed by taxes, no production, in neither $G$ nor GDP; the taxes of 55 do not enter either.}
\passo{d) bridge from 1{,}100}{Start from Prova~01 Q1(c) and add each change, by \eqref{q1:G1}:}
\begin{align*}
\text{GDP} &= (1{,}100)\underbrace{{}+(80)}_{\text{the app, new }C}\ \underbrace{{}+(30)}_{\text{government services}}\
\underbrace{{}-(100)}_{\text{A's imported raw materials}}\\
&= 1{,}110 ,
\end{align*}
which is \eqref{q8:Y}; the 100 of raw materials is no longer domestic value added. Selling 150 of B's goods abroad instead of at
home moves 150 from $C$ to $X$ in \eqref{q1:G4} and leaves GDP unchanged; the dividends, benefits and taxes change nothing.
\intuicao{Imports cancel twice over: the 100 of raw materials is inside the price of B's goods ($C$ and $X$), so it must come out
through $-M$. The software is counted exactly as in Prova~01 Q1(c): as investment, never as the consumption of the workers who
built it --- what they buy with their wages is already in $C$.}
\refbib{Kurlat (2020), \S1.1, Ex.~1.6 (imports and exports), p.~19; Exs.~1.7--1.8 (government at cost, transfers), p.~20; Jones
(2020), ch.~2, p.~22 (intellectual property products in $I$). GDP vs.\ GNI: Williamson (2014), solutions manual, ch.~2, problem~2(c).}
\rubrica{
\item a) 6 points: each approach 1.8, agreement 0.6. Only one approach: max.\ 2.
\item Software counted as consumption: max.\ 3 for a) and 0 for its line in c). Imports not subtracted: max.\ 3 for a).
\item Transfers or taxes added to $G$ or to income: max.\ 3 for a).
\item b) 2 points; GNI $=$ GDP $-40$ (all dividends treated as paid abroad): 1. c) 0.5 per classification. d) 2 points.
}
\end{sol}

% -----------------------------------------------------------------------------
\questao{(11 points) \emph{Capitalise or expense? Depreciation and measured growth.} Two years of the Prova~01 economy. A's
and B's value added is all paid out as factor income (wages and profits; the split is not needed). All values in reais.
\begin{itemize}[leftmargin=1.4em]
\item \textbf{Year 1} is Prova~01 Q1(c): A sells inputs to B for 600; B sells final goods to households for 1{,}000; N sells
advertising to B for 100 and uses that 100 to pay workers who build software for N's own use. The software is ready at the end of
year 1, will be used for four years and depreciates by $100/4=25$ a year from year 2 on.
\item \textbf{Year 2}: A, B and the advertising are as in year 1, and no new software is built. N launches ad-free premium
accounts: households buy accounts for 50, and B buys accounts for 20 for its marketing staff to use at work. N pays its support
staff wages of 30 and has no other costs.
\end{itemize}}
\noindent\textbf{a)} (3 points) Year 1: compute GDP by the three approaches (i) when the software is capitalised (current convention) and
(ii) when it is expensed as a current cost of N (the older convention).\\
\textbf{b)} (3 points) Year 2, software capitalised: compute GDP by the three approaches, and net domestic product (NDP).\\
\textbf{c)} (2 points) Compute GDP growth from year 1 to year 2 under each convention, and explain in words why the sign differs.\\
\textbf{d)} (1.5 points) In words: why do the households' 50 of premium accounts enter $C$ while B's 20 do not?\\
\textbf{e)} (1.5 points) In words: a student wrote that the 100 of software is ``consumption, because the workers spend their
wages''. Explain what is wrong, using the three columns.
\resolespaco
\begin{sol}
\begin{notacao}
\item A, B, N; $Q_j$, $Z_j$, $VA_j$, $W_j$, $\Pi_j$ (operating surplus, gross of depreciation) and the identities
\eqref{q1:G1}--\eqref{q1:G4}: as in Question~1. A superscript $1$ or $2$ marks the year.
\item $C$, $I\ge0$: consumption and investment ($G=X=M=0$). $D_N\ge0$: depreciation of N's software.
\item NDP: net domestic product. $g$: GDP growth rate from year 1 to year 2 (any sign).
\end{notacao}

\passo{a)(i) capitalised}{A's value added is its sales, 600. By \eqref{q1:G2}, B's output minus A's inputs and the advertising:}
\begin{align}
VA_B &= (1{,}000)-(600)-(100) \nonumber\\
&= 300. \label{q9:vaB}
\end{align}
N's output is the ads plus the software, valued at cost:
\begin{align}
VA_N &= (100)+(100) \nonumber\\
&= 200. \label{q9:vaN}
\end{align}
Summing with \eqref{q1:G1}:
\begin{align}
\text{GDP}^1_{\mathrm{cap}} &= (600)+(300)+(200) \nonumber\\
&= 1{,}100. \label{q9:Y1c}
\end{align}
Income: N pays wages of 100, so by \eqref{q1:G3} and \eqref{q9:vaN} its operating surplus is
\begin{align*}
\Pi_N &= (200)-(100)\\
&= 100 .
\end{align*}
Adding A's and B's factor income, N's wages and N's surplus, by \eqref{q1:G3}:
\begin{align*}
\text{GDP} &= (600)+(300)+(100)+(100)\\
&= 1{,}100 .
\end{align*}
Expenditure, by \eqref{q1:G4} with $C=1{,}000$ and $I=100$:
\begin{align*}
\text{GDP} &= (1{,}000)+(100)\\
&= 1{,}100 .
\end{align*}
\passo{a)(ii) expensed}{The software is not output, only a cost, so $VA_N=100$ (the ads). By \eqref{q1:G1}, with $VA_A=600$ and
\eqref{q9:vaB}:}
\begin{align}
\text{GDP}^1_{\mathrm{exp}} &= (600)+(300)+(100) \nonumber\\
&= 1{,}000. \label{q9:Y1e}
\end{align}
Income: N's profit is the ads minus the wages, by \eqref{q1:G3}:
\begin{align*}
\Pi_N &= (100)-(100)\\
&= 0 .
\end{align*}
Adding up incomes, by \eqref{q1:G3}:
\begin{align*}
\text{GDP} &= (600)+(300)+(100)+(0)\\
&= 1{,}000 .
\end{align*}
Expenditure: $C=1{,}000$ and $I=0$, so by \eqref{q1:G4}:
\begin{align*}
\text{GDP} &= (1{,}000)+(0)\\
&= 1{,}000 .
\end{align*}
\begin{center}\footnotesize\begin{tabular}{@{}lrr@{\qquad}lrr@{\qquad}lrr@{}}\toprule
Production & (i) & (ii) & Income & (i) & (ii) & Expenditure & (i) & (ii)\\\midrule
A & 600 & 600 & A & 600 & 600 & Consumption & 1{,}000 & 1{,}000\\
B & 300 & 300 & B & 300 & 300 & Investment & 100 & 0\\
N & 200 & 100 & N: wages & 100 & 100 & & & \\
 & & & N: surplus & 100 & 0 & & & \\\midrule
GDP & 1{,}100 & 1{,}000 & GDP & 1{,}100 & 1{,}000 & GDP & 1{,}100 & 1{,}000\\\bottomrule\end{tabular}\end{center}
\passo{b) year 2}{By \eqref{q1:G2}, B's output minus A's inputs, the advertising and the 20 of accounts:}
\begin{align}
VA^2_B &= (1{,}000)-(600)-(100)-(20) \nonumber\\
&= 280. \label{q9:vaB2}
\end{align}
N has no inputs; its output is the ads plus all premium accounts:
\begin{align}
VA^2_N &= (100)+(50)+(20) \nonumber\\
&= 170. \label{q9:vaN2}
\end{align}
Summing $VA_A=600$, \eqref{q9:vaB2} and \eqref{q9:vaN2} with \eqref{q1:G1}:
\begin{align}
\text{GDP}^2 &= (600)+(280)+(170) \nonumber\\
&= 1{,}050. \label{q9:Y2}
\end{align}
Income: N pays wages of 30, so by \eqref{q1:G3} and \eqref{q9:vaN2} its gross operating surplus is
\begin{align*}
\Pi^2_N &= (170)-(30)\\
&= 140 .
\end{align*}
The software depreciates by a quarter of its cost each year, and the surplus splits into that depreciation and net profit:
\begin{align*}
D_N &= \frac{(100)}{4}\\
&= 25,\\
\Pi^2_N-D_N &= (140)-(25)\\
&= 115 .
\end{align*}
Adding up incomes, by \eqref{q1:G3}:
\begin{align*}
\text{GDP}^2 &= (600)+(280)+(30)+(140)\\
&= 1{,}050 .
\end{align*}
Expenditure: no investment, no exports; by \eqref{q1:G4}:
\begin{align*}
\text{GDP}^2=C &= (1{,}000)+(50)\\
&= 1{,}050 .
\end{align*}
NDP subtracts depreciation from GDP:
\begin{equation}
\text{NDP}=\text{GDP}-D_N. \label{q9:ndp}
\end{equation}
Substitute \eqref{q9:Y2} and $D_N=25$ into \eqref{q9:ndp}:
\begin{align*}
\text{NDP} &= (1{,}050)-(25)\\
&= 1{,}025 .
\end{align*}
\passo{c) growth}{Capitalised: dividing \eqref{q9:Y2} by \eqref{q9:Y1c},}
\begin{align*}
g_{\mathrm{cap}} &= \frac{(1{,}050)}{(1{,}100)}-1\\
&= -4.55\% .
\end{align*}
Expensed: year-2 GDP is also \eqref{q9:Y2} (no software is built in year 2); dividing it by \eqref{q9:Y1e},
\begin{align*}
g_{\mathrm{exp}} &= \frac{(1{,}050)}{(1{,}000)}-1\\
&= +5.00\% .
\end{align*}
Under capitalisation year 1 contains a final good, the software, that year 2 does not produce, so
measured GDP falls although consumption rises from 1{,}000 to 1{,}050; under expensing the software never appears, and GDP only
sees the rise in consumption. The convention changes the level of GDP in the year the asset is built and, through it, the growth rate.
\passo{d) intermediate vs final}{The households use the accounts up themselves: final consumption, $C$ in \eqref{q1:G4}. B uses the
same accounts as an input for producing and selling its goods: intermediate consumption, in $Z_B$ of \eqref{q1:G2}, so B's value
added falls from \eqref{q9:vaB} to \eqref{q9:vaB2}:}
\begin{align*}
VA_B-VA^2_B &= (300)-(280)\\
&= 20 .
\end{align*}
That 20 is already inside the price of B's goods; adding it to $C$ would count it twice. N's value added rises by the full 70 in
both cases ($50+20$ in \eqref{q9:vaN2}).
\passo{e) the exam error}{The wages are on the \emph{income} side, \eqref{q1:G3}: N's workers earn 100, and what they buy with it
is part of $C=1{,}000$, which the statement keeps unchanged. The expenditure side, \eqref{q1:G4}, asks who uses the \emph{output}:
the software is kept by N and used for four years, so it is investment. Calling it consumption counts the workers' income as a final use and, with $C$ fixed at 1{,}000, leaves
the expenditure column at 1{,}000 against 1{,}100 in the other two --- the three columns stop agreeing, which flags the error.}
\intuicao{Capitalising treats own-account software like a machine bought from outside: an asset produced now that pays off later, its
wear recorded as depreciation (in NDP, not in GDP). Expensing treats it as a cost, like electricity. Same economy, different GDP.}
\refbib{Kurlat (2020), \S1.1, Ex.~1.3 (investment), p.~17; Ex.~1.10 (depreciation; GDP is gross), p.~21. Jones (2020), ch.~2,
p.~22 (intellectual property products --- software, R\&D --- in $I$) and p.~25 (effect on the labor share).}
\rubrica{
\item a) 3 points: 1.5 per convention, the three approaches each. Only GDP totals: max.\ 1.
\item b) 3 points: three approaches 2.4, NDP 0.6. Depreciation subtracted from GDP: max.\ 1.5 for b).
\item c) 2 points: both rates 1, explanation 1. d) 1.5 points. e) 1.5 points; ``it is consumption'' repeated: 0 for e).
}
\end{sol}
